\section{Coincident products: an explicit holomorphic completion}
\label{cub:sec}

The newest coincident-point report evaluates all quadratic Stieltjes
moments and explicitly asks for the cubic case, beginning with
$\FP\int_0^1\psi(x)^3\,dx$ \cite{IncomingCoincident}.
At three factors the local subtraction has a new feature: a pair of
singular factors and the linear term of the remaining factor together
produce a logarithmic divergence.  We give its exact spectral
counterterm.  The construction works at every derivative order and,
locally, at every number of factors.

\subsection{A finite local rule for any number of factors}

Use the lowering spectral convention
\begin{equation}\label{cub:generator}
 g(\lambda,x)=\zeta(1-\lambda,x)+\frac1\lambda
   =\sum_{m\ge0}\gamma_m(x)\frac{\lambda^m}{m!}.
\end{equation}
The pole at $\lambda=0$ is removable at each $x>0$.
For $r\ge0$ put
\[
 c_r(\lambda)=(-1)^r(1-\lambda)_r,
 \qquad h_r(\lambda,x)=\partial_x^r g(\lambda,1+x).
\]
The Hurwitz shift relation gives the exact splitting
\begin{equation}\label{cub:local-split}
 \partial_x^r g(\lambda,x)
 =c_r(\lambda)x^{\lambda-r-1}+h_r(\lambda,x).
\end{equation}
Every $h_r$ is holomorphic jointly in $\lambda$ near zero and $x$ near
zero.  Its required Taylor coefficients are explicit:
\begin{equation}\label{cub:h-coefficients}
 [x^k]h_r(\lambda,x)=
 \begin{cases}
 \zeta(1-\lambda)+1/\lambda,&r=k=0,\\[2pt]
 c_{r+k}(\lambda)\zeta(1+r+k-\lambda)/k!,&r+k>0.
 \end{cases}
\end{equation}

Let $d\ge2$ and $\boldsymbol r=(r_1,\ldots,r_d)$ be nonnegative
integers.  For $\varnothing\ne S\subseteq\{1,\ldots,d\}$ write
\[
 \lambda_S=\sum_{i\in S}\lambda_i,
 \qquad \nu_S=\sum_{i\in S}(r_i+1)-1,
\]
and define the finite numerator
\begin{equation}\label{cub:subset-numerator}
 A_S(\boldsymbol\lambda)=
 \left(\prod_{i\in S}c_{r_i}(\lambda_i)\right)
 [x^{\nu_S}]\prod_{j\notin S}h_{r_j}(\lambda_j,x).
\end{equation}
The empty smooth product is $1$.  In particular, $A_{\{1,\ldots,d\}}=0$
because $\nu_{\{1,\ldots,d\}}\ge1$.
Let $\mathcal L_{\boldsymbol r}$ denote the meromorphic continuation
of the ordinary integral of $\prod_i\partial_x^{r_i}g(\lambda_i,x)$.
It is initially defined, for example, with every
$\Re\lambda_i>r_i+1$, away from the removable parameter expressions.

\begin{theorem}[Subset completion in the unit coordinate]
\label{cub:general-completion}
The expression
\begin{equation}\label{cub:completion}
 \mathcal F_{\boldsymbol r}(\boldsymbol\lambda)
 =\mathcal L_{\boldsymbol r}(\boldsymbol\lambda)
   -\sum_{\varnothing\ne S\subsetneq\{1,\ldots,d\}}
            \frac{A_S(\boldsymbol\lambda)}{\lambda_S}
\end{equation}
is holomorphic in a neighborhood of the spectral origin.  Its Taylor
coefficients are exactly
\begin{equation}\label{cub:all-moments}
 \FP\int_0^1\prod_{i=1}^d\gamma_{m_i}^{(r_i)}(x)\,dx
 =\left(\prod_i m_i!\right)
 [\lambda_1^{m_1}\cdots\lambda_d^{m_d}]
          \mathcal F_{\boldsymbol r}(\boldsymbol\lambda).
\end{equation}
Here $\FP$ is the constant term of the cutoff integral in the fixed
coordinate $x$, with upper endpoint $1$ and no additional finite term.
The complete numerators in \eqref{cub:completion} are subtracted;
restricting a numerator to its polar hyperplane generally changes the
finite answer.
\end{theorem}

\begin{proof}
Fix $0<b<1$.  Expanding the exact product in
\eqref{cub:local-split} gives one term for each singular subset $S$,
namely
\[
 \left(\prod_{i\in S}c_{r_i}(\lambda_i)\right)
 x^{\lambda_S-\nu_S-1}
 \prod_{j\notin S}h_{r_j}(\lambda_j,x).
\]
For each nonempty $S$, subtract the Taylor polynomial of its smooth
product through degree $\nu_S$.  The remainder is normally integrable
on $(0,b)$ for all parameters in a sufficiently small common polydisc;
fixed parameter derivatives only add powers of $\log x$.
The Taylor monomial of degree $k$ has integral proportional to
\[
 \frac{b^{\lambda_S+k-\nu_S}}{\lambda_S+k-\nu_S}.
\]
Only $k=\nu_S$ has a denominator vanishing at the spectral origin.
Its meromorphic integral is $A_S b^{\lambda_S}/\lambda_S$.
The generating function of its coordinate finite parts is
$A_S(b^{\lambda_S}-1)/\lambda_S$.  Their difference is precisely
$A_S/\lambda_S$.  All other denominators are nonzero near the origin
and their Taylor coefficients equal the corresponding cutoff constant
terms.  The integral on $[b,1]$ is holomorphic and has the same
coefficient-extraction property.  Summing the finitely many subsets
proves both assertions.  Changing $b$ redistributes ordinary integrals
and Taylor integrals equally on both sides, so the result does not
depend on this auxiliary split.
\end{proof}

This theorem is a local completion rule, not a claim that an arbitrary
depth integral reduces to single zeta values.  For three factors its
global coordinates are the classical Tornheim function and ordinary
zeta functions.  That explicit global formula comes next.

\subsection{The cubic global kernel and all its argument derivatives}

Define, initially in its domain of absolute convergence,
\begin{equation}\label{cub:T-def}
 T(A,B,C)=\sum_{m,n\ge1}\frac{1}{m^A n^B(m+n)^C}.
\end{equation}
We always use its meromorphic continuation, with the order of any
restriction specified.  The classical Hurwitz--Tornheim connection is
due to earlier work, including Espinosa and Moll \cite{EM2006}; the
contribution here is its complete coordinate normalization at the
coincident Stieltjes point.

Put $\Lambda=\lambda_1+\lambda_2+\lambda_3$.  If $\{i,j,k\}$ is
$\{1,2,3\}$, define
\begin{align}
 K_2(a,b)&=\frac{2\Gamma(a)\Gamma(b)}{(2\pi)^{a+b}}
       \cos\frac{\pi(a-b)}2\,\zeta(a+b),\label{cub:K2}\\
 K_3(\boldsymbol\lambda)
 &=\frac{2\prod_i\Gamma(\lambda_i)}{(2\pi)^\Lambda}
   \sum_{k=1}^3
    \cos\frac{\pi(\lambda_i+\lambda_j-\lambda_k)}2
       T(\lambda_i,\lambda_j,\lambda_k).
       \label{cub:K3}
\end{align}
In \eqref{cub:K3}, the two indices other than $k$ may be taken in
increasing order, since $T(A,B,C)=T(B,A,C)$.

\begin{proposition}[Finite spectral expression]
\label{cub:global}
The kernels $K_2$ and $K_3$ are the meromorphic continuations of the
integrals of two and three factors $\zeta(1-\lambda_i,x)$,
respectively.  For $d=3$, the meromorphic integral in
Theorem~\ref{cub:general-completion} is explicitly
\begin{align}
 \mathcal L_{\boldsymbol r}(\boldsymbol\lambda)
 =\left(\prod_i c_{r_i}(\lambda_i)\right)
 \bigg\{&K_3(\lambda_1-r_1,\lambda_2-r_2,\lambda_3-r_3)
       \notag\\
 &+\sum_{k:r_k=0}\frac{
       K_2(\lambda_i-r_i,\lambda_j-r_j)}{\lambda_k}
       +\frac{\mathbf1_{\boldsymbol r=\boldsymbol0}}
                    {\lambda_1\lambda_2\lambda_3}\bigg\}.
 \label{cub:Lr}
\end{align}
Together with \eqref{cub:subset-numerator}, this is a finite
Tornheim--zeta coefficient formula for every cubic Stieltjes moment,
at arbitrary Stieltjes and argument-derivative orders.
\end{proposition}

\begin{proof}
For $\Re\lambda_i>1$, the nonzero Fourier coefficient of
$\zeta(1-\lambda_i,x)$ at $n$ is
\[
 \Gamma(\lambda_i)(2\pi|n|)^{-\lambda_i}
              e^{-\pi i\lambda_i\operatorname{sgn}(n)/2}.
\]
The Fourier series are absolutely convergent.  Integrating their
product keeps $n_1+n_2+n_3=0$, with all $n_i\ne0$.
There are three choices of the negative frequency and their three
sign reversals.  On each positive cone the negative magnitude is
$m+n$, giving exactly \eqref{cub:K3}.  The two-factor calculation
gives \eqref{cub:K2}; it is the classical Hurwitz product formula
\cite{EM2002}.  Finite endpoint Taylor subtraction supplies
meromorphic continuation of the original integrals, and uniqueness
extends these identities.  Finally,
$\partial_x^r g=c_r\zeta(1+r-\lambda,x)$ for $r>0$, whereas
$g=\zeta(1-\lambda,x)+1/\lambda$ for $r=0$.
Expand the three factors.  Integrals with just one Hurwitz factor
vanish by meromorphic continuation of
$\int_0^1\zeta(s,x)\,dx=0$.  The other terms give \eqref{cub:Lr}.
\end{proof}

For undifferentiated factors set
$h(\lambda)=\zeta(1-\lambda)+1/\lambda$.  Formula
\eqref{cub:completion} becomes particularly short:
\begin{align}
 \mathcal F_{\boldsymbol0}={}&K_3+
   \sum_{k=1}^3\frac{K_2(\lambda_i,\lambda_j)}{\lambda_k}
   +\frac{1}{\lambda_1\lambda_2\lambda_3}
   -\sum_{i=1}^3\frac{h(\lambda_j)h(\lambda_k)}{\lambda_i}
       \notag\\
 &+\sum_{k=1}^3\frac{(1-\lambda_k)\zeta(2-\lambda_k)}
                              {\lambda_i+\lambda_j}.
 \label{cub:Fexplicit}
\end{align}
The last line is necessary.  It is the pair-subset correction absent
from the separated-singularity triple formula.  Neither it nor the
singleton correction can be suppressed before restricting spectral
directions.

\subsection{The stated cubic digamma target}

Let
\[
 \omega(t)=T(t,t,t),\quad L=\log(2\pi),\quad
 E(t)=\Gamma(1+t)e^{-Lt},\quad
 \gamma=\gamma_0(1).
\]
The symbol $\omega^{(j)}(0)$ means a derivative of this specified
one-variable continuation.  It does not mean a derivative of a
jointly holomorphic $T$ at $(0,0,0)$.

\begin{theorem}[Diagonal second derivative and cubic digamma moment]
\label{cub:diagonal}
The diagonal continuation is holomorphic at zero and satisfies
\begin{equation}\label{cub:omega-jets}
 \omega(0)=\frac13,\qquad \omega'(0)=L,
 \qquad \boxed{\omega''(0)=L^2-4\zeta''(0)}.
\end{equation}
The first two values are classical \cite{BorweinDilcher}.
The coordinate-normalized cubic moment is
\begin{align}
 \boxed{\displaystyle\FP\int_0^1\psi(x)^3\,dx}
 ={}&-\omega'''(0)-8\zeta'''(0)
       -12(\gamma+L)\zeta''(0)-L^3-6\gamma L^2
       \notag\\
 &+5\gamma^3+6\gamma\gamma_1
       -\frac32\gamma\zeta(2)
       +\frac32\bigl(\zeta(2)+\zeta'(2)\bigr).
 \label{cub:psi-cube}
\end{align}
Thus the requested cubic finite part is exactly reduced to one third
derivative of the diagonal Tornheim function and ordinary zeta and
Stieltjes constants.  No reduction of $\omega'''(0)$ to ordinary zeta
constants is asserted.
\end{theorem}

\begin{proof}
On the diagonal, \eqref{cub:Fexplicit} reads
\begin{align}
 F(t):=\mathcal F_{\boldsymbol0}(t,t,t)
 ={}&\frac{6E(t)^3\cos(\pi t/2)\omega(t)
              +6E(t)^2\zeta(2t)+1}{t^3}
       \notag\\
 &-\frac{3h(t)^2}{t}
       +\frac{3(1-t)\zeta(2-t)}{2t}.
 \label{cub:Fdiagonal}
\end{align}
The left side is holomorphic by
Theorem~\ref{cub:general-completion}.  Solving this identity for
$\omega(t)$ proves its holomorphy at zero, since $E(0)=1$ and
$\cos0=1$.  Write
\[
 E(t)=\exp\left[-(\gamma+L)t+
                 \frac{\zeta(2)}2t^2-\frac{\zeta(3)}3t^3+O(t^4)\right].
\]
Cancellation of the coefficients of $t^{-3}$ and $t^{-2}$ gives the
first two values of \eqref{cub:omega-jets}, using
$\zeta(0)=-1/2$ and $\zeta'(0)=-L/2$.
The numerator coefficient at $t^2$ is
\[
 3\omega''(0)+3\gamma^2-3L^2-\tfrac32\zeta(2)+12\zeta''(0).
\]
The other two terms of \eqref{cub:Fdiagonal} contribute
$-3\gamma^2+3\zeta(2)/2$ to the coefficient of $t^{-1}$.
This proves the second-derivative formula.

The numerator coefficient at $t^3$, after those substitutions, is
\[
 \omega'''(0)+8\zeta'''(0)+12(\gamma+L)\zeta''(0)
 +L^3+6\gamma L^2-5\gamma^3+\tfrac32\gamma\zeta(2).
\]
The remaining constant terms are
$-6\gamma\gamma_1-3(\zeta(2)+\zeta'(2))/2$.
Finally $F(0)=\FP\int\gamma_0(x)^3dx=-\FP\int\psi(x)^3dx$,
which proves \eqref{cub:psi-cube}.
\end{proof}

A direct, convergent integral gives an independent way to evaluate
the same left side.  The expansion
$\psi(x)=-x^{-1}-\gamma+\zeta(2)x-\zeta(3)x^2+\cdots$ gives
\begin{align}
 \FP\int_0^1\psi(x)^3\,dx
 ={}&\int_0^1\left[\psi(x)^3+\frac1{x^3}
       +\frac{3\gamma}{x^2}
       -\frac{3\zeta(2)-3\gamma^2}{x}\right]dx
       +\frac12+3\gamma.
 \label{cub:direct-cube}
\end{align}
The integrand has a finite limit at zero.  Independent truncated
Mellin-series evaluation of $\omega'''(0)$ and local-series quadrature
of \eqref{cub:direct-cube} agree to more than thirty decimal places:
\begin{align}
 \omega''(0)&=11.403217934867761113621595961662\ldots,\notag\\
 \omega'''(0)&=86.657866601100073652903276858762\ldots,\notag\\
 \FP\int_0^1\psi(x)^3dx
   &=1.966262734391534340664375338411\ldots.
 \label{cub:numerics}
\end{align}
These are numerical diagnostics accompanying the analytic proof;
the implementation does not claim interval-certified final digits.

\begin{corollary}[A cubic derivative specialization]
\label{cub:cubic-derivative}
In the same coordinate convention,
\begin{equation}\label{cub:psi-square-prime}
 \FP\int_0^1\psi(x)^2\psi'(x)\,dx
             =\zeta(3)-2\gamma\zeta(2).
\end{equation}
\end{corollary}
\begin{proof}
Take the constant term of
$\int_\epsilon^1\psi^2\psi'=(\psi(1)^3-\psi(\epsilon)^3)/3$.
The constant coefficient of $\psi(\epsilon)^3$ is
$-\gamma^3+6\gamma\zeta(2)-3\zeta(3)$.
\end{proof}

More generally, a finite-part integral of a derivative of a product
equals its upper-endpoint value minus the constant term of its local
expansion.  This supplies a terminating family of relations between
the cubic moments in \eqref{cub:all-moments}; it does not impose a
spurious independent value on the remaining primitive cubic moment.

\subsection{Why a spectral direction must be recorded}

The reduction at fixed $A=B=0$ gives
\begin{equation}\label{cub:C-axis}
 T(0,0,C)=\zeta(C-1)-\zeta(C).
\end{equation}
Indeed, $m+n=N$ has $N-1$ positive solutions.  Continuing this
one-variable identity gives $5/12$ at $C=0$, whereas
\eqref{cub:omega-jets} gives $1/3$ on the diagonal.
This direction dependence is already explicit in
Borwein and Dilcher \cite[Theorem~2.7]{BorweinDilcher}.
It is not an inconsistency.  The multivariable function is meromorphic
at intersecting polar hyperplanes, and its restrictions carry different
finite information.  The completed function
$\mathcal F_{\boldsymbol r}$ is genuinely holomorphic; the uncompleted
$T$ generally is not.  Section~\ref{sp:sec:rays} computes the first two
derivatives along an arbitrary transverse ray.
